% GENERATED FILE. Edit canonical lessons, then run npm run generate:books.
% canonical-source-hash: fnv1a64:e6c6daa779574713
% canonical-lessons: GE-C490-schwiegervater, GE-C490-stiefvater, GE-C490-urgrossvater, GE-C490-zwilling, GE-C490-einzelkind

\chapter{Schwiegervater und Stiefvater}
\label{ch:ge-schwiegervater-und-stiefvater}
\hlchaptermodality{490}

\begin{chapteropening}
\textbf{By the end of this chapter:} \emph{I can talk about in-laws and step-families.}

Complete the last lesson of chapter 490: I can talk about in-laws and step-families.
\end{chapteropening}

\section[der Schwiegervater]{der Schwiegervater — father-in-law}

\label{lesson:GE-C490-schwiegervater}



\begin{quote}
Before the new one: say the German for niece, then the German for brother-in-law.
\end{quote}

\subsection*{You'll want to know: der Schwiegervater}
\textbf{der Schwiegervater} — \textquotedblleft{}father-in-law\textquotedblright{}.

\begin{quote}
Mein Schwiegervater kocht am Sonntag für alle.
\end{quote}

\begin{quote}
\emph{My father-in-law cooks for everyone on Sunday.}
\end{quote}

Plural: \emph{die Schwiegerväter}. The same first part gives \emph{die Schwiegermutter} (mother-in-law) and \emph{die Schwiegereltern} (parents-in-law).

\begin{cousinweb}
\emph{Schwiegervater} is built on \emph{Schwieger}, inherited from Old High German \emph{swigar} \textquotedblleft{}mother-in-law\textquotedblright{}, from Proto-Germanic \emph{swegrō} (EtymDB), and \emph{Vater} \textquotedblleft{}father\textquotedblright{}.
\end{cousinweb}

\begin{grammarlens}[title={What you've built}]
One more word for telling, explaining and describing at B1.
\end{grammarlens}

\subsection*{Your turn}
\begin{itemize}
\raggedright
  \item \emph{Say it:} \emph{der Schwiegervater}
  \item \emph{Say it:} \emph{der Schwiegervater}, once more
  \item \emph{Say it:} a sentence of your own with \emph{der Schwiegervater}
  \item \emph{Recall:} say the German for niece, then the German for brother-in-law, then say \emph{der Schwiegervater} again
\end{itemize}

\begin{tcolorbox}[breakable,colback=teal!4,colframe=teal!35!black,title={Before you move on}]
What does der Schwiegervater mean? (\textquotedblleft{}Father-in-law\textquotedblright{}.) Say it once more.
\end{tcolorbox}
\section[der Stiefvater]{der Stiefvater — stepfather}

\label{lesson:GE-C490-stiefvater}



\begin{quote}
Before the new one: say the German for brother-in-law, then the German for father-in-law.
\end{quote}

\subsection*{You'll want to know: der Stiefvater}
\textbf{der Stiefvater} — \textquotedblleft{}stepfather\textquotedblright{}.

\begin{quote}
Sie versteht sich gut mit ihrem Stiefvater.
\end{quote}

\begin{quote}
\emph{She gets on well with her stepfather.}
\end{quote}

Plural: \emph{die Stiefväter}. Likewise \emph{die Stiefmutter}, \emph{das Stiefkind}.

\begin{cousinweb}
\emph{Stiefvater} is a Germanic compound: EtymDB sets it beside English \emph{stepfather}, from Proto-Germanic \emph{steupafadēr}. The first part is usually said to have meant \textquotedblleft{}bereaved\textquotedblright{}.
\end{cousinweb}

\begin{grammarlens}[title={What you've built}]
One more word for telling, explaining and describing at B1.
\end{grammarlens}

\subsection*{Your turn}
\begin{itemize}
\raggedright
  \item \emph{Say it:} \emph{der Stiefvater}
  \item \emph{Say it:} \emph{der Stiefvater}, once more
  \item \emph{Say it:} a sentence of your own with \emph{der Stiefvater}
  \item \emph{Recall:} say the German for brother-in-law, then the German for father-in-law, then say \emph{der Stiefvater} again
\end{itemize}

\begin{tcolorbox}[breakable,colback=teal!4,colframe=teal!35!black,title={Before you move on}]
What does der Stiefvater mean? (\textquotedblleft{}Stepfather\textquotedblright{}.) Say it once more.
\end{tcolorbox}
\section[der Urgroßvater]{der Urgroßvater — great-grandfather}

\label{lesson:GE-C490-urgrossvater}



\begin{quote}
Before the new one: say the German for father-in-law, then the German for stepfather.
\end{quote}

\subsection*{You'll want to know: der Urgroßvater}
\textbf{der Urgroßvater} — \textquotedblleft{}great-grandfather\textquotedblright{}.

\begin{quote}
Mein Urgroßvater war Bauer.
\end{quote}

\begin{quote}
\emph{My great-grandfather was a farmer.}
\end{quote}

Plural: \emph{die Urgroßväter}. \emph{Ur-} adds one more generation: \emph{die Urgroßmutter}, \emph{der Urenkel}.

\begin{cousinweb}
\emph{Urgroßvater} is built on \emph{ur-} \textquotedblleft{}original\textquotedblright{} and \emph{Großvater}; \emph{groß} is inherited from Old High German \emph{grōz}, from Proto-Germanic \emph{grautaz} (EtymDB).
\end{cousinweb}

\begin{grammarlens}[title={What you've built}]
One more word for telling, explaining and describing at B1.
\end{grammarlens}

\subsection*{Your turn}
\begin{itemize}
\raggedright
  \item \emph{Say it:} \emph{der Urgroßvater}
  \item \emph{Say it:} \emph{der Urgroßvater}, once more
  \item \emph{Say it:} a sentence of your own with \emph{der Urgroßvater}
  \item \emph{Recall:} say the German for father-in-law, then the German for stepfather, then say \emph{der Urgroßvater} again
\end{itemize}

\begin{tcolorbox}[breakable,colback=teal!4,colframe=teal!35!black,title={Before you move on}]
What does der Urgroßvater mean? (\textquotedblleft{}Great-grandfather\textquotedblright{}.) Say it once more.
\end{tcolorbox}
\section[der Zwilling]{der Zwilling — twin}

\label{lesson:GE-C490-zwilling}



\begin{quote}
Before the new one: say the German for stepfather, then the German for great-grandfather.
\end{quote}

\subsection*{You'll want to know: der Zwilling}
\textbf{der Zwilling} — \textquotedblleft{}twin\textquotedblright{}.

\begin{quote}
Meine Schwester hat Zwillinge bekommen.
\end{quote}

\begin{quote}
\emph{My sister has had twins.}
\end{quote}

Plural: \emph{die Zwillinge}. \emph{Eineiige Zwillinge} — identical twins.

\begin{cousinweb}
\emph{Zwilling} is inherited from Middle High German \emph{zwillinc} (EtymDB), built on the old stem \emph{zwi-} \textquotedblleft{}two\textquotedblright{}; English \emph{twin} is related.
\end{cousinweb}

\begin{grammarlens}[title={What you've built}]
One more word for telling, explaining and describing at B1.
\end{grammarlens}

\subsection*{Your turn}
\begin{itemize}
\raggedright
  \item \emph{Say it:} \emph{der Zwilling}
  \item \emph{Say it:} \emph{der Zwilling}, once more
  \item \emph{Say it:} a sentence of your own with \emph{der Zwilling}
  \item \emph{Recall:} say the German for stepfather, then the German for great-grandfather, then say \emph{der Zwilling} again
\end{itemize}

\begin{tcolorbox}[breakable,colback=teal!4,colframe=teal!35!black,title={Before you move on}]
What does der Zwilling mean? (\textquotedblleft{}Twin\textquotedblright{}.) Say it once more.
\end{tcolorbox}
\section[das Einzelkind]{das Einzelkind — only child}

\label{lesson:GE-C490-einzelkind}



\begin{quote}
Before the new one: say the German for great-grandfather, then the German for twin.
\end{quote}

\subsection*{You'll want to know: das Einzelkind}
\textbf{das Einzelkind} — \textquotedblleft{}only child\textquotedblright{}.

\begin{quote}
Als Einzelkind hatte ich mein eigenes Zimmer.
\end{quote}

\begin{quote}
\emph{As an only child I had my own room.}
\end{quote}

Plural: \emph{die Einzelkinder}.

\begin{cousinweb}
\emph{Einzelkind} is built on \emph{einzeln} \textquotedblleft{}single\textquotedblright{} and \emph{Kind} \textquotedblleft{}child\textquotedblright{}, inherited from Middle High German \emph{kint} and Old High German \emph{kind} (EtymDB).
\end{cousinweb}

\begin{grammarlens}[title={What you've built}]
One more word for telling, explaining and describing at B1.
\end{grammarlens}

\subsection*{Your turn}
\begin{itemize}
\raggedright
  \item \emph{Say it:} \emph{das Einzelkind}
  \item \emph{Say it:} \emph{das Einzelkind}, once more
  \item \emph{Say it:} a sentence of your own with \emph{das Einzelkind}
  \item \emph{Recall:} say the German for great-grandfather, then the German for twin, then say \emph{das Einzelkind} again
\end{itemize}

\begin{tcolorbox}[breakable,colback=teal!4,colframe=teal!35!black,title={Before you move on}]
What does das Einzelkind mean? (\textquotedblleft{}Only child\textquotedblright{}.) Say it once more.
\end{tcolorbox}
